\section{Experto técnico II: modelo del mundo, retroalimentación por refuerzo y ciclo cerrado del tráfico}

La sección anterior explicó qué es VLA; esta sección explica cómo se entrena y se valida. Li Xiang (李想) mencionó que Li Auto (理想) usa un modelo del mundo para generar datos de entrenamiento: el modelo conduce del punto A al punto B, y el resultado se evalúa con retroalimentación sobre comodidad, colisiones, reglas de tránsito y otros aspectos. Aquí se aprecia un patrón muy claro: la flota real aporta los datos, el modelo del mundo genera escenarios y aplica los exámenes, VLA se actualiza con esos datos y esa retroalimentación, y al final se despliega de nuevo en los vehículos.

\begin{figure}[H]
\centering
\includegraphics[width=0.96\textwidth]{world-model-loop.png}
\caption{Ciclo cerrado del modelo del mundo del tráfico: los datos de conducción reales, la generación por simulación y el entrenamiento de VLA se impulsan mutuamente. Diagrama conceptual propio, elaborado a partir de la conversación en 00:47:30--00:48:10 y 01:10:06--01:10:48.}
\end{figure}

\begin{knowledgebox}{Lectura de la figura: el modelo del mundo cumple tres funciones}
Primero, puede servir como sistema de examen para evaluar el desempeño de VLA en distintos escenarios de tráfico. Segundo, puede generar datos de entrenamiento que cubren escenarios difíciles de recopilar con frecuencia en las vías reales. Tercero, en el futuro podría formar parte del sistema de operación de vehículos autónomos y ayudar a los vehículos sin conductor a entender y predecir el mundo del tráfico.
\end{knowledgebox}

\subsection{El modelo del mundo no es un «generador de imágenes»}

En el contexto de la conducción autónoma, el modelo del mundo no se limita a generar videos atractivos. Debe ser sensible a las consecuencias de las acciones: cómo evoluciona el mundo del tráfico si el vehículo cambia de carril, acelera, esquiva o espera. También debe ser sensible a la seguridad y a las reglas: las colisiones, los frenados bruscos, pasarse un semáforo en rojo, pisar las líneas del carril, colarse en la fila y la comodidad no son solo cuestiones visuales, sino cuestiones del resultado de la tarea.

\begin{importantbox}{La fórmula mínima del modelo del mundo}
El modelo del mundo del tráfico puede escribirse, de forma aproximada, como:
\[
s_{t+1} = f(s_t, a_t, c_t).
\]
Aquí, \(s_t\) representa el estado actual del tráfico, \(a_t\) la acción del vehículo, \(c_t\) las restricciones y reglas del escenario, y \(s_{t+1}\) el estado en el instante siguiente. Lo verdaderamente difícil es lograr que \(f\) sea lo bastante confiable en los escenarios de cola larga de las vías reales.
\end{importantbox}

\subsection{Costo de validación y escala del modelo}

Li Xiang mencionó que el costo de validación bajó de 180,000 yuanes a 4,000 yuanes por cada 10,000 kilómetros. Esto indica que la ruta de VLA y el modelo del mundo no es solo una cuestión de capacidad del modelo, sino que también cambia la economía de la validación. Si la conducción autónoma solo pudiera validarse acumulando kilometraje en vías reales, el costo sería altísimo; el modelo del mundo y la simulación pueden reducir ciertos costos de prueba, pero al final la validación en el mundo real sigue siendo indispensable como respaldo.

\begin{warningbox}{La simulación no puede sustituir directamente al mundo real}
El modelo del mundo puede reducir los costos de exploración y validación, pero no puede reemplazar por completo al mundo real. Los comportamientos de cola larga que la simulación no cubre, el ruido de los sensores, las obras viales, el clima extremo y las acciones irracionales de las personas pueden seguir apareciendo en las vías reales.
\end{warningbox}

\subsection{L3, L4 y el límite de cómputo a bordo del vehículo}

Ya se explicó cómo el modelo del mundo reduce los costos de entrenamiento y validación; esta subsección lleva el problema a los niveles de conducción autónoma y a las restricciones del cómputo a bordo. En la entrevista, Li Xiang relacionó la capacidad L3/L4 con la escala de los modelos en la nube y a bordo. A su juicio, la generación actual de cómputo corresponde aproximadamente a la capacidad L3, y L4 todavía requiere más capacidades y más condiciones. El valor didáctico de esta apreciación es que el nivel de conducción autónoma no es un simple interruptor de software, sino que lo determinan en conjunto la escala del modelo, el cómputo a bordo, el sistema de validación, la regulación y la responsabilidad operativa.

\begin{knowledgebox}{Términos clave: L3 y L4}
L3 suele significar que, en condiciones específicas, el sistema puede asumir la tarea de conducción, pero una persona debe tomar el control cuando el sistema lo solicita; L4 significa que, dentro de un dominio de diseño operativo delimitado, el sistema puede funcionar sin depender de que una persona tome el control. La diferencia entre ambos no está solo en las puntuaciones del modelo, sino en los límites de responsabilidad, el diseño de redundancia, la regulación y el sistema de operación.
\end{knowledgebox}

\subsection{Resumen de la sección}

El modelo del mundo lleva a VLA de ser «un modelo que sabe conducir» a ser «un sistema que puede entrenarse, examinarse y operarse». Reduce el costo de validación, pero no elimina la incertidumbre del mundo real; eleva el techo de desempeño, pero también exige más cómputo, más datos y un sistema de seguridad más sólido.
